\documentclass[11pt]{ctexart}
\usepackage[a4paper,margin=1in]{geometry}
\usepackage{longtable,booktabs}
\title{Geant4-Agent 回归测试报告（中文）}
\author{自动化评测流程}
\date{2026-03-07\_154815}
\begin{document}
\maketitle
\section{统一 Pass/Fail 总表}
\begin{longtable}{p{0.06\linewidth}p{0.08\linewidth}p{0.07\linewidth}p{0.08\linewidth}p{0.07\linewidth}p{0.07\linewidth}p{0.07\linewidth}p{0.07\linewidth}p{0.09\linewidth}p{0.06\linewidth}p{0.21\linewidth}}
\toprule
ID & 领域 & 风格 & 参数完整性 & 期望初始 & 实际初始 & 期望最终 & 实际最终 & 缺参(初->终) & 结果 & 失败原因 \\
\midrule
MTM01 & medical & fuzzy & missing & 失败 & 失败 & 通过 & 通过 & 13->0 & 通过 & - \\
MTM02 & aerospace & precise & missing & 失败 & 失败 & 通过 & 通过 & 3->0 & 通过 & - \\
MTM03 & industrial\_ndt & precise & missing & 失败 & 失败 & 通过 & 通过 & 11->0 & 通过 & - \\
MTM04 & physics & precise & missing & 失败 & 失败 & 通过 & 通过 & 8->0 & 通过 & - \\
MTM05 & nuclear\_engineering & fuzzy & missing & 失败 & 失败 & 通过 & 通过 & 5->0 & 通过 & - \\
MTM06 & security\_screening & fuzzy & missing & 失败 & 失败 & 通过 & 通过 & 12->0 & 通过 & - \\
MTM07 & semiconductor & precise & missing & 失败 & 失败 & 通过 & 通过 & 3->0 & 通过 & - \\
MTM08 & education & fuzzy & missing & 失败 & 失败 & 通过 & 通过 & 11->0 & 通过 & - \\
MTM09 & environmental\_radiation & fuzzy & missing & 失败 & 失败 & 通过 & 通过 & 3->0 & 通过 & - \\
MTM10 & high\_energy\_physics & precise & missing & 失败 & 失败 & 通过 & 通过 & 10->0 & 通过 & - \\
MTM11 & materials\_update & precise & full & 通过 & 通过 & 通过 & 通过 & 0->0 & 通过 & - \\
MTM12 & overwrite\_reject & precise & full & 通过 & 通过 & 通过 & 通过 & 0->0 & 通过 & - \\
MTM13 & medical\_detector & fuzzy & missing & 失败 & 失败 & 通过 & 通过 & 11->0 & 失败 & expected\_final\_mismatch:geometry.structure \\
MTM14 & industrial\_panel & precise & missing & 失败 & 失败 & 通过 & 通过 & 15->0 & 失败 & expected\_final\_mismatch:geometry.structure \\
MTM15 & shielding\_stack & precise & missing & 失败 & 失败 & 通过 & 失败 & 15->6 & 失败 & final\_complete\_mismatch, expected\_final\_mismatch:geometry.structure \\
MTM16 & nested\_target & precise & missing & 失败 & 失败 & 通过 & 通过 & 6->0 & 失败 & expected\_final\_mismatch:geometry.structure \\
MTM17 & coaxial\_shield & fuzzy & missing & 失败 & 失败 & 通过 & 通过 & 14->0 & 失败 & expected\_final\_mismatch:geometry.structure \\
MTM18 & boolean\_cavity & precise & missing & 失败 & 失败 & 通过 & 通过 & 12->0 & 失败 & expected\_final\_mismatch:geometry.structure \\
MTM19 & ring\_complete & precise & full & 通过 & 通过 & 通过 & 通过 & 0->0 & 失败 & expected\_final\_mismatch:geometry.structure \\
MTM20 & grid\_complete & precise & full & 通过 & 通过 & 通过 & 通过 & 0->0 & 失败 & expected\_final\_mismatch:geometry.structure, expected\_final\_mismatch:materials.selected\_materials.0 \\
MTM21 & layered\_detector & fuzzy & missing & 失败 & 失败 & 通过 & 通过 & 15->0 & 失败 & expected\_final\_mismatch:geometry.structure \\
MTM22 & embedded\_sensor & precise & full & 通过 & 通过 & 通过 & 通过 & 0->0 & 失败 & expected\_final\_mismatch:geometry.structure \\
MTM23 & layered\_shield\_complete & precise & full & 通过 & 失败 & 通过 & 失败 & 6->6 & 失败 & initial\_complete\_mismatch, final\_complete\_mismatch, expected\_final\_mismatch:geometry.structure \\
MTM24 & boolean\_complete & precise & full & 通过 & 通过 & 通过 & 通过 & 0->0 & 失败 & expected\_final\_mismatch:geometry.structure \\
\bottomrule
\end{longtable}
\section{基准对比}
\begin{itemize}
\item pass\_rate=0.500
\item initial\_alignment\_rate=0.958
\item final\_alignment\_rate=0.917
\item avg\_turns\_per\_case=2.250
\item missing\_reduction\_rate=1.000
\item confirm\_apply\_success\_rate=1.000
\item rollback\_after\_confirm\_rate=0.000
\item internal\_leak\_turn\_rate=0.037
\item repeated\_assistant\_case\_rate=0.000
\item stale\_confirm\_case\_rate=0.000
\item delta\_pass\_rate=-0.450
\item regressions=无
\item improvements=无
\end{itemize}
\appendix
\section{附录：完整对话与内部处理链（自然语言）}
\subsection{案例 MTM01（medical）}
\textbf{对话：}
\begin{itemize}
\item 用户：我先要一个用于医学成像的初步辐照场景，细节还没定，你先告诉我还缺什么。
\item 系统：你提到的医学成像场景，具体需要什么样的几何形状呢？比如是简单的盒体、圆柱，还是更复杂的结构？另外，如果选择盒体，它在x方向上的尺寸大概是多少？
\item 用户：采用 300 mm x 300 mm x 300 mm 的水箱，材料 G4\_WATER；源用束流 gamma，能量 6 MeV，位置 (0,0,-400) mm，方向 +z；物理列表 QBBC；输出 root。
\item 系统：好的，配置已完成！目前我们已确认的设置包括：体积与材料绑定、输出路径和输出格式。
\end{itemize}
\textbf{内部链路（自然语言）：}
\begin{itemize}
\item 第 1 轮：动作=ask\_clarification，阶段=geometry，闭环状态=False。
\item 缺失字段：geometry.structure、geometry.params.module\_x、geometry.params.module\_y、geometry.params.module\_z、materials.volume\_material\_map、source.type、source.particle、source.energy、source.position、source.direction、physics.physics\_list、output.format、output.path。
\item LLM 使用=True；回退=无；后端=llm\_semantic\_frame+runtime\_semantic。
\item 第 2 轮：动作=finalize，阶段=materials，闭环状态=True。
\item 缺失字段：无。
\item LLM 使用=True；回退=无；后端=llm\_slot\_frame+runtime\_semantic。
\end{itemize}
\subsection{案例 MTM02（aerospace）}
\textbf{对话：}
\begin{itemize}
\item 用户：I need a first-pass proton irradiation setup for a silicon device; keep it simple for now.
\item 系统：Geometry: updated geometry structure, box size along x, box size along y. Materials: updated selected materials, volume-to-material binding. Source: updated source type, particle type, source energy. Physics: still needs physics list. Output: still needs output format, output path.
\item 用户：Use a silicon box 50 mm x 50 mm x 5 mm, beam source proton 100 MeV from (0,0,-100) mm along +z, physics FTFP\_BERT, output json.
\item 系统：This would overwrite an already confirmed value.

Pending changes:
- Box size along x: 1000.0 → 50.0
- Box size along y: 1000.0 → 50.0

Reply **'confirm'** to apply these changes, or **'keep existing'** to keep the current values.
\item 用户：confirm
\item 系统：Configuration complete. The following items are confirmed: box size along x, box size along y, box size along z, and output path.
\end{itemize}
\textbf{内部链路（自然语言）：}
\begin{itemize}
\item 第 1 轮：动作=summarize\_progress，阶段=materials，闭环状态=False。
\item 缺失字段：physics.physics\_list、output.format、output.path。
\item LLM 使用=True；回退=无；后端=llm\_slot\_frame+runtime\_semantic。
\item 第 2 轮：动作=confirm\_overwrite，阶段=source，闭环状态=False。
\item 缺失字段：无。
\item LLM 使用=True；回退=无；后端=llm\_slot\_frame+runtime\_semantic。
\item 第 3 轮：动作=finalize，阶段=physics，闭环状态=True。
\item 缺失字段：无。
\item LLM 使用=True；回退=无；后端=llm\_semantic\_frame+runtime\_semantic+confirmed\_pending\_overwrite。
\end{itemize}
\subsection{案例 MTM03（industrial\_ndt）}
\textbf{对话：}
\begin{itemize}
\item 用户：我要做一个 gamma 穿透铁靶的小测试，但我还没把完整参数说完。
\item 系统：好的，我们已确认了材料（铁）、粒子类型（gamma）和物理列表。

接下来还需要你补充几何部分的信息，包括结构类型、盒体尺寸，以及将体积与材料绑定的设置。
\item 用户：请设为半径 40 mm、半长 80 mm 的铁圆柱，材料 G4\_Fe；点源 gamma，能量 2 MeV，位置 (0,0,-200) mm，方向 +z；物理列表 Shielding；输出 csv。
\item 系统：配置已完成。当前已确认的设置包括：体积与材料绑定、输出路径、物理列表和输出格式。
\end{itemize}
\textbf{内部链路（自然语言）：}
\begin{itemize}
\item 第 1 轮：动作=summarize\_progress，阶段=geometry，闭环状态=False。
\item 缺失字段：geometry.structure、geometry.params.module\_x、geometry.params.module\_y、geometry.params.module\_z、materials.volume\_material\_map、source.type、source.energy、source.position、source.direction、output.format、output.path。
\item LLM 使用=True；回退=无；后端=llm\_slot\_frame+runtime\_semantic。
\item 第 2 轮：动作=finalize，阶段=materials，闭环状态=True。
\item 缺失字段：无。
\item LLM 使用=True；回退=无；后端=llm\_slot\_frame+runtime\_semantic。
\end{itemize}
\subsection{案例 MTM04（physics）}
\textbf{对话：}
\begin{itemize}
\item 用户：Gamma point source 1 MeV at (0,0,-100) toward +z; please build the rest.
\item 系统：Geometry: still needs geometry structure, box size along x. Materials: still needs volume-to-material binding. Source: updated source type, particle type, source energy. Physics: still needs physics list. Output: still needs output format, output path.
\item 用户：Set copper box 1 m x 1 m x 1 m, physics FTFP\_BERT, output hdf5.
\item 系统：Your configuration is complete. The following items have been confirmed: volume-to-material binding, output path, output format, and physics list.
\end{itemize}
\textbf{内部链路（自然语言）：}
\begin{itemize}
\item 第 1 轮：动作=summarize\_progress，阶段=geometry，闭环状态=False。
\item 缺失字段：geometry.structure、geometry.params.module\_x、geometry.params.module\_y、geometry.params.module\_z、materials.volume\_material\_map、physics.physics\_list、output.format、output.path。
\item LLM 使用=True；回退=无；后端=llm\_slot\_frame+runtime\_semantic。
\item 第 2 轮：动作=finalize，阶段=materials，闭环状态=True。
\item 缺失字段：无。
\item LLM 使用=True；回退=无；后端=llm\_slot\_frame+runtime\_semantic。
\end{itemize}
\subsection{案例 MTM05（nuclear\_engineering）}
\textbf{对话：}
\begin{itemize}
\item 用户：我想做一个中子屏蔽的概念验证，现在先给一个缺参开头。
\item 系统：好的，当前配置进度如下：

**已设置完成：**
- **几何**：结构已确定为盒体。
- **材料**：材料已选定并绑定到体积。
- **源**：已确定使用中子源，并设置了能量。
- **物理**：物理列表已选定。
- **输出**：输出格式和文件路径已设置。

**仍需补充以下参数才能继续：**
- **几何尺寸**：盒体在x、y、z三个方向的具体尺寸。
- **源信息**：中子的发射位置和方向。

请提供这些信息，我们就能生成完整的模拟配置了。
\item 用户：使用 200 mm x 200 mm x 200 mm 的铅块，材料 G4\_Pb；点源 neutron，能量 14 MeV，位置 (0,0,-50) mm，方向 +z；物理列表 QGSP\_BERT\_HP；输出 xml。
\item 系统：我注意到这可能会覆盖已经确认的设置。具体来说：

*   **已选材料**：将从 `G4\_Cu` 改为 `G4\_Pb`。
*   **粒子类型**：将从 `gamma` 改为 `neutron`。

如果您确认应用这些修改，请直接回复“确认”；如果您希望保留原来的值，请直接回复“保持原值”。
\item 用户：确认
\item 系统：好的，配置已完成。目前我们已确认的设置包括：体积与材料绑定、输出路径、输出格式和物理列表。
\end{itemize}
\textbf{内部链路（自然语言）：}
\begin{itemize}
\item 第 1 轮：动作=summarize\_progress，阶段=geometry，闭环状态=False。
\item 缺失字段：geometry.params.module\_x、geometry.params.module\_y、geometry.params.module\_z、source.position、source.direction。
\item LLM 使用=True；回退=无；后端=llm\_slot\_frame+runtime\_semantic。
\item 第 2 轮：动作=confirm\_overwrite，阶段=materials，闭环状态=False。
\item 缺失字段：无。
\item LLM 使用=True；回退=无；后端=llm\_slot\_frame+runtime\_semantic。
\item 第 3 轮：动作=finalize，阶段=source，闭环状态=True。
\item 缺失字段：无。
\item LLM 使用=True；回退=无；后端=llm\_semantic\_frame+runtime\_semantic+confirmed\_pending\_overwrite。
\end{itemize}
\subsection{案例 MTM06（security\_screening）}
\textbf{对话：}
\begin{itemize}
\item 用户：Need a quick cargo-screening toy setup with forward photons, but I have not fixed the target shape, material, source energy, or output format yet.
\item 系统：Here's a summary of your setup so far:

**Completed:**
*   **Source:** Particle type is set to photons.

**Still Needed:**
*   **Geometry:** Structure and box dimensions (X, Y, Z).
*   **Materials:** Assignment of a material to your geometry volume.
*   **Source:** Type, energy, and position.
*   **Physics:** Selection of a physics list.
*   **Output:** Format and file path.
\item 用户：Use steel cylinder target radius 60 mm half-length 120 mm, beam gamma 3 MeV from (0,0,-500) mm to +z, physics Shielding, output root.
\item 系统：Your configuration is complete. The following items are now confirmed: the material is bound to your volume, the output path and format are set, and the Shielding physics list is active.
\end{itemize}
\textbf{内部链路（自然语言）：}
\begin{itemize}
\item 第 1 轮：动作=summarize\_progress，阶段=geometry，闭环状态=False。
\item 缺失字段：geometry.structure、geometry.params.module\_x、geometry.params.module\_y、geometry.params.module\_z、materials.volume\_material\_map、source.type、source.energy、source.position、source.direction、physics.physics\_list、output.format、output.path。
\item LLM 使用=True；回退=无；后端=llm\_slot\_frame+runtime\_semantic。
\item 第 2 轮：动作=finalize，阶段=materials，闭环状态=True。
\item 缺失字段：无。
\item LLM 使用=True；回退=无；后端=llm\_slot\_frame+runtime\_semantic。
\end{itemize}
\subsection{案例 MTM07（semiconductor）}
\textbf{对话：}
\begin{itemize}
\item 用户：我想搭一个硅薄片的中子照射测试，但先只给一个简述。
\item 系统：目前我们已经完成了几何结构（盒体尺寸）、材料（硅）和粒子源（中子能量与位置）的设置。

接下来还需要你补充以下信息：
1. 物理列表（例如：QGSP\_BERT）
2. 输出格式（例如：ROOT、CSV）
3. 输出文件路径

请提供这些信息，我们就可以继续生成配置了。
\item 用户：使用 20 mm x 20 mm x 1 mm 的硅片，材料 G4\_Si；各向同性 neutron，能量 2 MeV，位置 (0,0,0) mm；物理列表 QGSP\_BERT\_HP；输出 json。
\item 系统：我注意到这次修改会覆盖已经确认的几何尺寸。具体来说：
- 盒体 x 方向尺寸：1000.0 → 20.0
- 盒体 y 方向尺寸：1000.0 → 20.0

请确认你的选择：
- 如果同意应用这些修改，请直接回复“确认”。
- 如果想保留原来的尺寸，请直接回复“保持原值”。
\item 用户：确认
\item 系统：配置已完成。当前已确认的参数包括：源位置、源能量、源类型以及盒体 z 方向尺寸。
\end{itemize}
\textbf{内部链路（自然语言）：}
\begin{itemize}
\item 第 1 轮：动作=summarize\_progress，阶段=materials，闭环状态=False。
\item 缺失字段：physics.physics\_list、output.format、output.path。
\item LLM 使用=True；回退=无；后端=llm\_slot\_frame+runtime\_semantic。
\item 第 2 轮：动作=confirm\_overwrite，阶段=source，闭环状态=False。
\item 缺失字段：无。
\item LLM 使用=True；回退=无；后端=llm\_slot\_frame+runtime\_semantic。
\item 第 3 轮：动作=finalize，阶段=physics，闭环状态=True。
\item 缺失字段：无。
\item LLM 使用=True；回退=无；后端=llm\_semantic\_frame+runtime\_semantic+confirmed\_pending\_overwrite。
\end{itemize}
\subsection{案例 MTM08（education）}
\textbf{对话：}
\begin{itemize}
\item 用户：For classroom demo, I want gamma attenuation in copper; you can ask only essential missing parts.
\item 系统：Geometry: still needs geometry structure, box size along x. Materials: updated selected materials; still needs volume-to-material binding. Source: updated particle type; still needs source type, source energy. Physics: updated physics list. Output: still needs output format, output path.
\item 用户：Okay, copper box 100 mm x 100 mm x 100 mm, point gamma 1 MeV at (0,0,-200) mm, direction +z, physics FTFP\_BERT, output csv.
\item 系统：You've already set the physics list to "Shielding." Changing it to "FTFP\_BERT" would overwrite that.

Please reply **'confirm'** to apply the change to FTFP\_BERT, or **'keep existing'** to keep the current Shielding list.
\item 用户：confirm
\item 系统：Configuration complete. The physics list, output path, output format, and source direction are now confirmed.
\end{itemize}
\textbf{内部链路（自然语言）：}
\begin{itemize}
\item 第 1 轮：动作=summarize\_progress，阶段=geometry，闭环状态=False。
\item 缺失字段：geometry.structure、geometry.params.module\_x、geometry.params.module\_y、geometry.params.module\_z、materials.volume\_material\_map、source.type、source.energy、source.position、source.direction、output.format、output.path。
\item LLM 使用=True；回退=无；后端=llm\_slot\_frame+runtime\_semantic。
\item 第 2 轮：动作=confirm\_overwrite，阶段=materials，闭环状态=False。
\item 缺失字段：无。
\item LLM 使用=True；回退=无；后端=llm\_slot\_frame+runtime\_semantic。
\item 第 3 轮：动作=finalize，阶段=source，闭环状态=True。
\item 缺失字段：无。
\item LLM 使用=True；回退=无；后端=llm\_semantic\_frame+runtime\_semantic+confirmed\_pending\_overwrite。
\end{itemize}
\subsection{案例 MTM09（environmental\_radiation）}
\textbf{对话：}
\begin{itemize}
\item 用户：我想先搭一个空气体积中的低能 gamma 传播场景，细节稍后补。
\item 系统：好的，目前我们已经完成了场景的基础搭建：

*   **几何**：已设定为长方体，并确定了X、Y方向的尺寸。
*   **材料**：已选定材料，并完成了与体积的绑定。
*   **源**：已设定为低能伽马源，并配置了粒子类型与能量。

接下来，还需要您补充以下信息才能完成配置：
1.  **物理列表**
2.  **输出格式**
3.  **输出路径**
\item 用户：设为 500 mm x 500 mm x 500 mm 的空气盒，材料 G4\_AIR；点源 gamma，能量 0.5 MeV，位置 (0,0,-300) mm，方向 +z；物理列表 QBBC；输出 root。
\item 系统：我注意到这会覆盖已经确认的几何尺寸。具体修改如下：

*   盒体 x 方向尺寸：1000.0 → 500.0
*   盒体 y 方向尺寸：1000.0 → 500.0

请确认您的选择：
*   如果要应用这些修改，请直接回复“确认”。
*   如果要保留原值，请直接回复“保持原值”。
\item 用户：确认
\item 系统：配置已完成。当前已确认的参数包括：源位置、源能量、源类型以及盒体的 Z 方向尺寸。
\end{itemize}
\textbf{内部链路（自然语言）：}
\begin{itemize}
\item 第 1 轮：动作=summarize\_progress，阶段=materials，闭环状态=False。
\item 缺失字段：physics.physics\_list、output.format、output.path。
\item LLM 使用=True；回退=无；后端=llm\_slot\_frame+runtime\_semantic。
\item 第 2 轮：动作=confirm\_overwrite，阶段=source，闭环状态=False。
\item 缺失字段：无。
\item LLM 使用=True；回退=无；后端=llm\_slot\_frame+runtime\_semantic。
\item 第 3 轮：动作=finalize，阶段=physics，闭环状态=True。
\item 缺失字段：无。
\item LLM 使用=True；回退=无；后端=llm\_semantic\_frame+runtime\_semantic+confirmed\_pending\_overwrite。
\end{itemize}
\subsection{案例 MTM10（high\_energy\_physics）}
\textbf{对话：}
\begin{itemize}
\item 用户：Set geometry as a sphere target; I have not decided radius or source details yet.
\item 系统：You've set the geometry to a sphere. Here's what we have and what's still needed:

**Completed:**
*   **Geometry:** Structure is a sphere. You've also provided box dimensions (likely for a world volume).

**Still Needed:**
*   **Geometry:** The sphere's radius.
*   **Materials:** Which material to assign to your volume.
*   **Source:** The source type, particle type, and its energy.
*   **Physics:** The physics list to use.
*   **Output:** The output format and file path.
\item 用户：Sphere target radius 80 mm, lead material, electron point source 5 MeV at (0,0,-200) mm toward +z, physics FTFP\_BERT, output hdf5.
\item 系统：Your configuration is complete. The following items have been confirmed: volume-to-material binding, output path, output format, and physics list.
\end{itemize}
\textbf{内部链路（自然语言）：}
\begin{itemize}
\item 第 1 轮：动作=summarize\_progress，阶段=geometry，闭环状态=False。
\item 缺失字段：geometry.params.child\_rmax、materials.volume\_material\_map、source.type、source.particle、source.energy、source.position、source.direction、physics.physics\_list、output.format、output.path。
\item LLM 使用=True；回退=无；后端=llm\_slot\_frame+runtime\_semantic。
\item 第 2 轮：动作=finalize，阶段=materials，闭环状态=True。
\item 缺失字段：无。
\item LLM 使用=True；回退=无；后端=llm\_slot\_frame+runtime\_semantic。
\end{itemize}
\subsection{案例 MTM11（materials\_update）}
\textbf{对话：}
\begin{itemize}
\item 用户：Create a complete setup: 100 mm x 100 mm x 100 mm copper box, point gamma 1 MeV at (0,0,-200) mm toward +z, physics QBBC, output json.
\item 系统：Your configuration is complete. The volume-to-material binding, output path, and output format have been confirmed.
\item 用户：Change the target material to G4\_Al and change the output format to root.
\item 系统：You're about to overwrite a confirmed value.

Pending changes:
- Selected materials: from G4\_Cu to G4\_Al
- Output format: from json to root

Please reply **'confirm'** to apply these changes, or **'keep existing'** to keep the current values.
\item 用户：confirm
\item 系统：Configuration complete. All items are confirmed: volume-to-material binding, output path, output format, and selected materials.
\end{itemize}
\textbf{内部链路（自然语言）：}
\begin{itemize}
\item 第 1 轮：动作=finalize，阶段=materials，闭环状态=True。
\item 缺失字段：无。
\item LLM 使用=True；回退=无；后端=llm\_slot\_frame+runtime\_semantic。
\item 第 2 轮：动作=confirm\_overwrite，阶段=source，闭环状态=False。
\item 缺失字段：无。
\item LLM 使用=True；回退=无；后端=llm\_slot\_frame+runtime\_semantic。
\item 第 3 轮：动作=finalize，阶段=physics，闭环状态=True。
\item 缺失字段：无。
\item LLM 使用=True；回退=无；后端=llm\_semantic\_frame+runtime\_semantic+confirmed\_pending\_overwrite。
\end{itemize}
\subsection{案例 MTM12（overwrite\_reject）}
\textbf{对话：}
\begin{itemize}
\item 用户：请创建一个完整配置：100 mm x 100 mm x 100 mm 的铜盒，材料 G4\_Cu；点源 gamma，能量 1 MeV，位置 (0,0,-200) mm，方向 +z；物理列表 QBBC；输出 json。
\item 系统：配置已完成。当前已确认的设置包括：体积与材料绑定、输出路径和输出格式。
\item 用户：把材料改成 G4\_Al，并把输出格式改成 root。
\item 系统：我注意到这次修改会覆盖已确认的配置项。具体变动如下：

- 已选材料：从 G4\_Cu 改为 G4\_Al
- 输出格式：从 json 改为 root

如果你确认应用这些修改，请直接回复“确认”；如果希望保留原来的值，请回复“保持原值”。
\item 用户：保持原值
\item 系统：好的，已保留您之前确认的设置，材料仍为 G4\_Cu，输出格式仍为 json。
\end{itemize}
\textbf{内部链路（自然语言）：}
\begin{itemize}
\item 第 1 轮：动作=finalize，阶段=materials，闭环状态=True。
\item 缺失字段：无。
\item LLM 使用=True；回退=无；后端=llm\_slot\_frame+runtime\_semantic。
\item 第 2 轮：动作=confirm\_overwrite，阶段=source，闭环状态=False。
\item 缺失字段：无。
\item LLM 使用=True；回退=无；后端=llm\_slot\_frame+runtime\_semantic。
\item 第 3 轮：动作=reject\_overwrite，阶段=physics，闭环状态=True。
\item 缺失字段：无。
\item LLM 使用=True；回退=无；后端=llm\_semantic\_frame+runtime\_semantic+rejected\_pending\_overwrite。
\end{itemize}
\subsection{案例 MTM13（medical\_detector）}
\textbf{对话：}
\begin{itemize}
\item 用户：I want a ring of detector modules around a small source region, but I have not fixed the ring radius or module size yet.
\item 系统：Here's where we stand:

**Geometry**
- Completed: You've chosen a torus-based structure for your detector ring and set the size of the central box.
- Still needed: The torus tube radius and its major radius (which defines the ring size).

**Materials**
- Still needed: You'll need to assign a material to your volumes.

**Source**
- Still needed: The type of source and the particle it emits.

**Physics \& Output**
- Still needed: A physics list, an output format, and a file path for your results.
\item 用户：Use a ring of 12 modules, each 8 mm x 10 mm x 2 mm, radius 40 mm, clearance 1 mm, material G4\_WATER; point gamma source 0.662 MeV at (0,0,0) mm; physics QBBC; output json.
\item 系统：This would overwrite an already confirmed value. Pending change: geometry structure: single\_torus -> single\_box; box size along x: None -> 8.0. Reply 'confirm' to apply it, or say 'keep existing' to keep the current value.
\item 用户：confirm
\item 系统：Your configuration is complete. The following items have been confirmed: the torus major radius, the torus tube radius, the volume-to-material binding, and the radius.
\end{itemize}
\textbf{内部链路（自然语言）：}
\begin{itemize}
\item 第 1 轮：动作=summarize\_progress，阶段=geometry，闭环状态=False。
\item 缺失字段：geometry.params.torus\_rmax、geometry.params.torus\_rtor、materials.volume\_material\_map、source.type、source.particle、source.energy、source.position、source.direction、physics.physics\_list、output.format、output.path。
\item LLM 使用=True；回退=无；后端=llm\_slot\_frame+runtime\_semantic。
\item 第 2 轮：动作=confirm\_overwrite，阶段=geometry，闭环状态=False。
\item 缺失字段：geometry.params.torus\_rmax、geometry.params.torus\_rtor。
\item LLM 使用=True；回退=无；后端=llm\_slot\_frame+runtime\_semantic。
\item 第 3 轮：动作=finalize，阶段=materials，闭环状态=True。
\item 缺失字段：无。
\item LLM 使用=True；回退=无；后端=llm\_semantic\_frame+runtime\_semantic+confirmed\_pending\_overwrite。
\end{itemize}
\subsection{案例 MTM14（industrial\_panel）}
\textbf{对话：}
\begin{itemize}
\item 用户：我想做一个二维阵列探测板，但 pitch 和材料还没定。
\item 系统：目前我们已经确定了几何结构类型和盒体在x、y、z三个方向的尺寸。接下来还需要补充体积与材料的绑定关系，以及源的设置，包括源类型、粒子类型、能量和位置。
\item 用户：设成 4 x 3 的盒体阵列，每个模块 10 mm x 10 mm x 2 mm，pitch\_x 12 mm，pitch\_y 14 mm，间隙 1 mm，材料 G4\_Si；点源 gamma 0.511 MeV，位置 (0,0,-80) mm，方向 +z；物理列表 QBBC；输出 root。
\item 系统：我注意到这将会覆盖已经确认的参数。具体修改如下：
- 盒体 x 方向尺寸：1000.0 → 10.0
- 盒体 y 方向尺寸：1000.0 → 10.0

请确认您的选择：如果应用这些修改，请回复“确认”；如果保留原值，请回复“保持原值”。
\item 用户：确认
\item 系统：好的，配置已经完成。目前我们确认的参数包括：y方向阵列数量、x方向阵列数量、盒体z方向尺寸、盒体y方向尺寸和盒体x方向尺寸。
\end{itemize}
\textbf{内部链路（自然语言）：}
\begin{itemize}
\item 第 1 轮：动作=summarize\_progress，阶段=materials，闭环状态=False。
\item 缺失字段：materials.volume\_material\_map、source.type、source.particle、source.energy、source.position、source.direction、physics.physics\_list、output.format、output.path、geometry.ask.grid.module\_size、geometry.ask.grid.count\_x、geometry.ask.grid.count\_y、geometry.ask.grid.pitch\_x、geometry.ask.grid.pitch\_y、geometry.ask.grid.clearance。
\item LLM 使用=True；回退=无；后端=llm\_slot\_frame+runtime\_semantic。
\item 第 2 轮：动作=confirm\_overwrite，阶段=source，闭环状态=False。
\item 缺失字段：geometry.ask.grid.count\_x、geometry.ask.grid.count\_y。
\item LLM 使用=True；回退=无；后端=llm\_slot\_frame+runtime\_semantic。
\item 第 3 轮：动作=finalize，阶段=physics，闭环状态=True。
\item 缺失字段：无。
\item LLM 使用=True；回退=无；后端=llm\_semantic\_frame+runtime\_semantic+confirmed\_pending\_overwrite。
\end{itemize}
\subsection{案例 MTM15（shielding\_stack）}
\textbf{对话：}
\begin{itemize}
\item 用户：Build a stacked absorber sandwich along z; I will provide the layer sizes next.
\item 系统：I've updated the geometry structure to a stacked absorber sandwich along the z-axis.

Here's what we have so far and what's still needed:

**Completed:**
*   **Geometry:** The structure is set. You've also provided the box sizes (X, Y, Z) for the world volume.

**Still Needed:**
*   **Geometry:** The radii for the three layers (r1, r2, r3).
*   **Materials:** Which material goes in each volume.
*   **Source:** The source type, particle type, and energy.
*   **Physics:** The physics list to use.
*   **Output:** The format and file path for saving results.

Please provide the three radii (r1, r2, r3) for the layers next.
\item 用户：Use stacked layers with footprint 20 mm x 20 mm, thicknesses 2 mm, 4 mm, 6 mm, layer clearance 0.5 mm, outer box 30 mm x 30 mm x 20 mm, container clearance 1 mm, material G4\_Al; beam gamma 1.25 MeV from (0,0,-120) mm along +z; physics Shielding; output csv.
\item 系统：Here's what we've set up this turn:

**Geometry:** Updated the outer box size (X, Y, Z). We still need the radii for the first two layers (r1, r2).
**Materials:** Selected material (G4\_Al) and assigned it to the volumes.
**Source:** Set the source type, particle (gamma), and energy (1.25 MeV).
**Physics:** Selected the Shielding physics list.
**Output:** Set the format to CSV and confirmed the output path.

To finish, please provide the two remaining layer radii: r1 and r2.
\end{itemize}
\textbf{内部链路（自然语言）：}
\begin{itemize}
\item 第 1 轮：动作=summarize\_progress，阶段=geometry，闭环状态=False。
\item 缺失字段：geometry.params.r1、geometry.params.r2、geometry.params.r3、geometry.params.z1、geometry.params.z2、geometry.params.z3、materials.volume\_material\_map、source.type、source.particle、source.energy、source.position、source.direction、physics.physics\_list、output.format、output.path。
\item LLM 使用=True；回退=无；后端=llm\_slot\_frame+runtime\_semantic。
\item 第 2 轮：动作=summarize\_progress，阶段=geometry，闭环状态=False。
\item 缺失字段：geometry.params.r1、geometry.params.r2、geometry.params.r3、geometry.params.z1、geometry.params.z2、geometry.params.z3。
\item LLM 使用=True；回退=无；后端=llm\_slot\_frame+runtime\_semantic。
\end{itemize}
\subsection{案例 MTM16（nested\_target）}
\textbf{对话：}
\begin{itemize}
\item 用户：我想在一个盒子里嵌一个圆柱靶，先别急着补全。
\item 系统：目前我们已经确定了几何结构类型和盒子的尺寸，也设好了源的位置、方向和类型。

接下来还需要补充这些信息：为体积指定材料、选择粒子类型和能量、设置物理列表，以及决定输出格式和路径。
\item 用户：外盒 120 mm x 120 mm x 120 mm，内嵌铜圆柱半径 20 mm、半长 40 mm，嵌套间隙 2 mm；材料 G4\_Cu；束流 gamma 1 MeV，位置 (0,0,-150) mm，方向 +z；物理列表 QBBC；输出 json。
\item 系统：我注意到新的尺寸设置会覆盖之前已经确认的数值。

具体修改如下：
- 盒体 x 方向尺寸：从 1000.0 改为 120.0
- 盒体 y 方向尺寸：从 1000.0 改为 120.0

如果你确定要应用这些修改，请回复“确认”；如果想保留原来的数值，请回复“保持原值”。
\item 用户：确认
\item 系统：配置已完成。当前已确认的参数包括：源位置、源类型以及盒体的 x、y、z 方向尺寸。
\end{itemize}
\textbf{内部链路（自然语言）：}
\begin{itemize}
\item 第 1 轮：动作=summarize\_progress，阶段=materials，闭环状态=False。
\item 缺失字段：materials.volume\_material\_map、source.particle、source.energy、physics.physics\_list、output.format、output.path。
\item LLM 使用=True；回退=无；后端=llm\_slot\_frame+runtime\_semantic。
\item 第 2 轮：动作=confirm\_overwrite，阶段=source，闭环状态=False。
\item 缺失字段：无。
\item LLM 使用=True；回退=无；后端=llm\_slot\_frame+runtime\_semantic。
\item 第 3 轮：动作=finalize，阶段=physics，闭环状态=True。
\item 缺失字段：无。
\item LLM 使用=True；回退=无；后端=llm\_semantic\_frame+runtime\_semantic+confirmed\_pending\_overwrite。
\end{itemize}
\subsection{案例 MTM17（coaxial\_shield）}
\textbf{对话：}
\begin{itemize}
\item 用户：I need a shell-style shielding toy model, but the shell thicknesses are still open.
\item 系统：Here's what we've set up so far:

**Geometry:** We have a shell-style structure. Still need the inner and outer radii (r1, r2) and the shell thicknesses (r3, z1, z2).
**Materials:** Need to assign materials to the volumes.
**Source:** Need the source type, particle type, and energy.
**Physics:** Physics list is configured.
**Output:** Need the output format and file path.

What would you like to do next?
\item 用户：Use a shell with inner radius 15 mm, thicknesses 5 mm and 8 mm, half length 40 mm, nested core radius 10 mm, nested core half length 30 mm, clearance 1 mm, material G4\_Pb; isotropic neutron source 2 MeV at (0,0,0) mm; physics Shielding; output hdf5.
\item 系统：Your configuration is complete. The following items are confirmed: volume-to-material binding, output path, physics list, and output format.
\end{itemize}
\textbf{内部链路（自然语言）：}
\begin{itemize}
\item 第 1 轮：动作=summarize\_progress，阶段=geometry，闭环状态=False。
\item 缺失字段：geometry.params.r1、geometry.params.r2、geometry.params.r3、geometry.params.z1、geometry.params.z2、geometry.params.z3、materials.volume\_material\_map、source.type、source.particle、source.energy、source.position、source.direction、output.format、output.path。
\item LLM 使用=True；回退=无；后端=llm\_slot\_frame+runtime\_semantic。
\item 第 2 轮：动作=finalize，阶段=materials，闭环状态=True。
\item 缺失字段：无。
\item LLM 使用=True；回退=无；后端=llm\_slot\_frame+runtime\_semantic。
\end{itemize}
\subsection{案例 MTM18（boolean\_cavity）}
\textbf{对话：}
\begin{itemize}
\item 用户：我想做一个挖空的盒体几何，但先只给你意图。
\item 系统：目前我们已经确定了使用挖空的盒体作为几何结构，接下来还需要补充一些具体信息：

- **几何**：需要提供盒体在 x、y、z 三个方向上的尺寸。
- **材料**：需要指定体积与材料的绑定关系。
- **源**：需要选择源的类型（例如粒子类型、能量等）。

请随时告诉我这些信息，我们可以继续配置。
\item 用户：做减法布尔体：A 为 40 mm x 40 mm x 40 mm 盒体，B 为 20 mm x 20 mm x 20 mm 盒体，从 A 中减去 B；材料 G4\_Al；点源 gamma 0.3 MeV，位置 (0,0,-60) mm，方向 +z；物理列表 QBBC；输出 xml。
\item 系统：配置已完成。当前已确认的设置包括：体积与材料绑定、输出路径、输出格式和物理列表。
\end{itemize}
\textbf{内部链路（自然语言）：}
\begin{itemize}
\item 第 1 轮：动作=summarize\_progress，阶段=geometry，闭环状态=False。
\item 缺失字段：geometry.params.module\_x、geometry.params.module\_y、geometry.params.module\_z、materials.volume\_material\_map、source.type、source.particle、source.energy、source.position、source.direction、physics.physics\_list、output.format、output.path。
\item LLM 使用=True；回退=无；后端=llm\_slot\_frame+runtime\_semantic。
\item 第 2 轮：动作=finalize，阶段=materials，闭环状态=True。
\item 缺失字段：无。
\item LLM 使用=True；回退=无；后端=llm\_slot\_frame+runtime\_semantic。
\end{itemize}
\subsection{案例 MTM19（ring\_complete）}
\textbf{对话：}
\begin{itemize}
\item 用户：Create a complete ring setup: 16 modules, each 6 mm x 6 mm x 2 mm, radius 32 mm, clearance 0.5 mm, material G4\_Si; point gamma source 0.511 MeV at (0,0,0) mm; physics QBBC; output json.
\item 系统：Your configuration is complete. The volume-to-material binding, output path, and module count have been confirmed.
\end{itemize}
\textbf{内部链路（自然语言）：}
\begin{itemize}
\item 第 1 轮：动作=finalize，阶段=materials，闭环状态=True。
\item 缺失字段：无。
\item LLM 使用=True；回退=无；后端=llm\_slot\_frame+runtime\_semantic。
\end{itemize}
\subsection{案例 MTM20（grid\_complete）}
\textbf{对话：}
\begin{itemize}
\item 用户：请创建完整的阵列配置：3 x 3 盒体阵列，每个模块 12 mm x 12 mm x 3 mm，pitch\_x 15 mm，pitch\_y 15 mm，间隙 1 mm，材料 G4\_CsI；点源 gamma 0.662 MeV，位置 (0,0,-90) mm，方向 +z；物理列表 QBBC；输出 root。
\item 系统：配置已完成。当前已确认的设置包括：体积与材料绑定、输出路径和输出格式。
\end{itemize}
\textbf{内部链路（自然语言）：}
\begin{itemize}
\item 第 1 轮：动作=finalize，阶段=materials，闭环状态=True。
\item 缺失字段：无。
\item LLM 使用=True；回退=无；后端=llm\_slot\_frame+runtime\_semantic。
\end{itemize}
\subsection{案例 MTM21（layered\_detector）}
\textbf{对话：}
\begin{itemize}
\item 用户：I want a stacked detector slab, but the individual layer thicknesses are not fixed yet.
\item 系统：Here's what we have so far:

**Geometry:** You've chosen a stacked detector slab structure and defined the overall box size. Next, we need the footprint (size) of each layer and their individual thicknesses.

**Materials:** We still need to assign materials to the volumes.

**Source:** We need the source type, particle type, and energy.

**Physics \& Output:** The physics list and output details are also still required.
\item 用户：Use stack footprint 25 mm x 25 mm, thicknesses 1 mm, 2 mm, 1 mm, layer clearance 0.2 mm, outer box 30 mm x 30 mm x 12 mm, container clearance 0.5 mm, material G4\_Si; beam electron 5 MeV from (0,0,-50) mm along +z; physics FTFP\_BERT; output json.
\item 系统：This would overwrite an already confirmed value. Pending change: geometry structure: stack -> single\_box; box size along x: None -> 30.0. Reply 'confirm' to apply it, or say 'keep existing' to keep the current value.
\item 用户：confirm
\item 系统：Configuration complete. The following items are now confirmed: overall box size (X, Y, Z) and the geometry structure.
\end{itemize}
\textbf{内部链路（自然语言）：}
\begin{itemize}
\item 第 1 轮：动作=summarize\_progress，阶段=geometry，闭环状态=False。
\item 缺失字段：geometry.graph\_program、materials.volume\_material\_map、source.type、source.particle、source.energy、source.position、source.direction、physics.physics\_list、output.format、output.path、geometry.ask.stack.footprint、geometry.ask.stack.thicknesses、geometry.ask.stack.layer\_clearance、geometry.ask.stack.parent\_size、geometry.ask.stack.nest\_clearance。
\item LLM 使用=True；回退=无；后端=llm\_semantic\_frame+runtime\_semantic。
\item 第 2 轮：动作=confirm\_overwrite，阶段=materials，闭环状态=False。
\item 缺失字段：无。
\item LLM 使用=True；回退=无；后端=llm\_slot\_frame+runtime\_semantic。
\item 第 3 轮：动作=finalize，阶段=source，闭环状态=True。
\item 缺失字段：无。
\item LLM 使用=True；回退=无；后端=llm\_semantic\_frame+runtime\_semantic+confirmed\_pending\_overwrite。
\end{itemize}
\subsection{案例 MTM22（embedded\_sensor）}
\textbf{对话：}
\begin{itemize}
\item 用户：请创建完整的嵌套配置：外盒 80 mm x 80 mm x 80 mm，内嵌铅圆柱半径 15 mm、半长 25 mm，嵌套间隙 1 mm；材料 G4\_Pb；点源 gamma 1 MeV，位置 (0,0,-100) mm，方向 +z；物理列表 Shielding；输出 csv。
\item 系统：配置已完成。当前已确认的设置包括：体积与材料绑定、输出路径和输出格式。
\end{itemize}
\textbf{内部链路（自然语言）：}
\begin{itemize}
\item 第 1 轮：动作=finalize，阶段=materials，闭环状态=True。
\item 缺失字段：无。
\item LLM 使用=True；回退=无；后端=llm\_slot\_frame+runtime\_semantic。
\end{itemize}
\subsection{案例 MTM23（layered\_shield\_complete）}
\textbf{对话：}
\begin{itemize}
\item 用户：Create a complete shell configuration: inner radius 12 mm, thicknesses 4 mm and 6 mm, half length 35 mm, nested core radius 8 mm, nested core half length 25 mm, clearance 1 mm, material G4\_Fe; isotropic gamma source 0.8 MeV at (0,0,0) mm; physics QBBC; output hdf5.
\item 系统：Here's a summary of your configuration progress.

**Completed this turn:**
*   **Geometry:** Structure, radius, and half length set.
*   **Materials:** Materials and their volume assignments set.
*   **Source:** Type, particle, and energy set.
*   **Physics:** Physics list set.
*   **Output:** Format and path set.

**Still needed for geometry:**
*   The inner radii (r1, r2, r3) and half lengths (z1, z2) for the nested shells.
\end{itemize}
\textbf{内部链路（自然语言）：}
\begin{itemize}
\item 第 1 轮：动作=summarize\_progress，阶段=geometry，闭环状态=False。
\item 缺失字段：geometry.params.r1、geometry.params.r2、geometry.params.r3、geometry.params.z1、geometry.params.z2、geometry.params.z3。
\item LLM 使用=True；回退=无；后端=llm\_slot\_frame+runtime\_semantic。
\end{itemize}
\subsection{案例 MTM24（boolean\_complete）}
\textbf{对话：}
\begin{itemize}
\item 用户：请创建完整布尔配置：做并集布尔体，A 为 30 mm x 30 mm x 20 mm 盒体，B 为 10 mm x 10 mm x 40 mm 盒体，与 A 做 union；材料 G4\_Al；束流 gamma 0.5 MeV，位置 (0,0,-70) mm，方向 +z；物理列表 QBBC；输出 json。
\item 系统：配置已完成！当前已确认的设置有：体积与材料绑定、输出路径，以及布尔体 B 的 x、y、z 尺寸。
\end{itemize}
\textbf{内部链路（自然语言）：}
\begin{itemize}
\item 第 1 轮：动作=finalize，阶段=materials，闭环状态=True。
\item 缺失字段：无。
\item LLM 使用=True；回退=无；后端=llm\_slot\_frame+runtime\_semantic。
\end{itemize}
\end{document}